% !TEX root = ../../ctfp-print.tex

\lettrine[lhang=0.17]{モ}{ナドが何のためにあるか}を理解した今――それは装飾された関数を合成するためだ――本当に興味深い問いは、なぜ装飾された関数が関数型プログラミングでそれほど重要なのかということだ。すでに一つの例として\code{Writer}モナドを見た。そこでは装飾によって、複数の関数呼び出しにわたってログを作成・蓄積することができた。他の方法では（グローバルな状態にアクセスして変更するなど）不純な関数を使って解決しなければならなかった問題が、純粋な関数で解決された。

\section{問題}

以下は、Eugenio Moggiの先駆的な論文から抜粋した類似問題の短いリストで、\footnote{Eugenio Moggi:
  \href{https://doi.org/10.1016/0890-5401(91)90052-4}{
    Notions of Computation and Monads}. Inf. Comput. 93(1): 55-92 (1991).}
いずれも伝統的に関数の純粋性を放棄することで解決されてきた。

\begin{itemize}
  \tightlist
  \item
        部分性（Partiality）：終了しないかもしれない計算
  \item
        非決定性（Nondeterminism）：多くの結果を返すかもしれない計算
  \item
        副作用（Side effects）：状態にアクセス・変更する計算

        \begin{itemize}
          \tightlist
          \item
                読み取り専用の状態、または環境
          \item
                書き込み専用の状態、またはログ
          \item
                読み書き可能な状態
        \end{itemize}
  \item
        例外（Exceptions）：失敗するかもしれない部分関数
  \item
        継続（Continuations）：プログラムの状態を保存して必要に応じて復元する能力
  \item
        対話型入力（Interactive Input）
  \item
        対話型出力（Interactive Output）
\end{itemize}

本当に驚くべきことは、これらすべての問題が同じ巧みなトリックを使って解決できることだ：装飾された関数に頼るのだ。もちろん、装飾はケースごとにまったく異なる。

この段階では、装飾がモナド的である必要はないことを認識しなければならない。合成を要求するとき――一つの装飾された関数をより小さな装飾された関数に分解できること――になって初めてモナドが必要になる。繰り返しになるが、それぞれの装飾が異なるため、モナド的合成は異なる方法で実装されるが、全体のパターンは同じだ。それは非常に単純なパターンだ：結合的で恒等元を備えた合成。

次のセクションはHaskellの例が豊富だ。圏論に戻りたい場合や、すでにHaskellのモナド実装に精通している場合は、読み飛ばしても構わない。

\section{解決策}

まず、\code{Writer}モナドをどのように使ったかを分析しよう。特定のタスクを実行する純粋な関数から始めた――引数を与えると特定の出力を生成する。この関数を別の関数に置き換えた。その関数は元の出力を文字列とペアにすることで装飾した。それがロギング問題への解決策だった。

そこで止まることはできなかった。一般に、モノリシックな解決策を扱いたくないからだ。ログを生成する一つの関数を、より小さなログを生成する関数に分解できる必要があった。それらの小さな関数の合成がモナドの概念へと導いた。

本当に驚くべきことは、関数の戻り値型を装飾するという同じパターンが、通常は純粋性を放棄することを必要とする多種多様な問題に機能することだ。リストを確認しながら、それぞれの問題に適用される装飾を特定しよう。

\subsection{部分性}

終了しないかもしれないすべての関数の戻り値型を「持ち上げられた」型に変更する――元の型のすべての値と特別な「ボトム」値$\bot$を含む型だ。たとえば、集合としての\code{Bool}型は二つの要素を持つ：\code{True}と\code{False}。持ち上げられた\code{Bool}は三つの要素を含む。持ち上げられた\code{Bool}を返す関数は\code{True}や\code{False}を生成するか、永遠に実行し続けるかもしれない。

面白いことに、Haskellのような遅延言語では、終わらない関数が実際に値を返すかもしれず、その値を次の関数に渡せる。この特別な値をボトムと呼ぶ。明示的に必要とされない限り（例えば、パターンマッチングの対象や出力として生成される場合）、プログラムの実行を停止させることなく受け渡しができる。すべてのHaskell関数は潜在的に非終了かもしれないため、Haskellの全ての型は持ち上げられていると仮定される。これが、単純な$\Set$ではなくHaskellの（持ち上げられた）型と関数の圏$\Hask$についてよく語る理由だ。ただし$\Hask$が真の圏かどうかは定かではない（この\urlref{http://math.andrej.com/2016/08/06/hask-is-not-a-category/}{Andrej Bauerの投稿}を参照）。

\subsection{非決定性}

関数が多くの異なる結果を返せる場合、すべてを一度に返してもよい。意味論的には、非決定的な関数は結果のリストを返す関数と等価だ。これは遅延ガーベジコレクション言語では非常に理にかなっている。たとえば、必要なのが一つの値だけなら、リストの先頭を取るだけでよく、末尾は評価されない。ランダムな値が必要なら、乱数生成器を使ってリストのn番目の要素を選べばよい。遅延評価によって無限のリストを返すことさえできる。

リストモナド――Haskellの非決定的計算の実装――では、\code{join}は\code{concat}として実装される。\code{join}はコンテナのコンテナを平坦化するはずだということを思い出してほしい――\code{concat}はリストのリストを一つのリストに連結する。\code{return}はシングルトンリストを作る：

\src{snippet01}
リストモナドのバインド演算子は、\code{fmap}に続いて\code{join}という一般的な公式で与えられ、この場合は次のようになる：

\src{snippet02}
ここで、自身がリストを生成する関数\code{k}がリスト\code{as}のすべての要素に適用される。結果はリストのリストで、\code{concat}を使って平坦化される。

プログラマーの観点からは、リストを扱うことは、非決定的な関数をループで呼び出したり、イテレーターを返す関数を実装したりするよりも簡単だ（ただし、\urlref{http://ericniebler.com/2014/04/27/range-comprehensions/}{モダンなC++}では、遅延レンジを返すことはHaskellでリストを返すこととほぼ等価になる）。

非決定性を創造的に使う良い例はゲームプログラミングにある。たとえば、コンピュータが人間とチェスを指す場合、対戦相手の次の手を予測できない。しかし、すべての可能な手のリストを生成して一つずつ分析することができる。同様に、非決定的なパーサーは与えられた式に対するすべての可能な解析のリストを生成するかもしれない。

リストを返す関数を非決定的と解釈できるにしても、リストモナドの応用はずっと広い。それは、リストを生成する計算をつなぎ合わせることが、命令型プログラミングで使われるループ――反復構造――の完璧な関数的代替となるからだ。単一のループはしばしば、リストの各要素にループ本体を適用する\code{fmap}を使って書き直せる。リストモナドでの\code{do}記法は複雑なネストされたループを置き換えるために使える。

私のお気に入りの例は、ピタゴラス数のトリプル――直角三角形の辺を形成できる正の整数のトリプル――を生成するプログラムだ。

\src{snippet03}
最初の行は、\code{z}が正の数の無限リスト\code{{[}1..{]}}から要素を取得することを示している。次に、\code{x}が1から\code{z}の間の数のリスト\code{{[}1..z{]}}から要素を取得する。最後に\code{y}が\code{x}から\code{z}の間の数のリストから要素を取得する。手元に三つの数$1 \leqslant x \leqslant y \leqslant z$がある。関数\code{guard}は\code{Bool}式を受け取り、ユニットのリストを返す：

\src{snippet04}
この関数（より大きなクラス\code{MonadPlus}のメンバー）は、ピタゴラス数でないトリプルを除外するために使われる。実際、バインド（または関連する演算子\code{>>}）の実装を見ると、空のリストを渡されると空のリストを生成することがわかる。一方、空でないリスト（ここでは、ユニット\code{{[}(){]}}を含むシングルトンリスト）を渡されると、バインドは継続、ここでは\code{return (x, y, z)}、を呼び出し、検証済みのピタゴラス数トリプルを含むシングルトンリストを生成する。これらすべてのシングルトンリストは、囲んでいるバインドによって連結され、最終的な（無限の）結果を生成する。もちろん、\code{triples}の呼び出し元はリスト全体を消費することはできないが、Haskellは遅延であるため問題ない。

通常なら三つのネストされたループのセットを必要とする問題が、リストモナドと\code{do}記法の助けを借りて劇的に単純化された。それだけでなく、Haskellではリスト内包表記を使ってさらにコードを単純化できる：

\src{snippet05}
これはリストモナド（厳密には\code{MonadPlus}）のさらなる糖衣構文だ。

他の関数型または命令型言語では、ジェネレーターやコルーチンという名目で同様の構造を見ることができる。

\subsection{読み取り専用の状態}

外部の状態または環境への読み取り専用アクセスを持つ関数は、常にその環境を追加引数として取る関数に置き換えられる。純粋な関数\code{(a, e) -> b}（\code{e}は環境の型）は、一見するとKleisli矢のように見えないかもしれない。しかしカリー化して\code{a -> (e -> b)}にすると、装飾が古い友人であるリーダー関手であると気づく：

\src{snippet06}
\code{Reader}を返す関数はミニ実行可能ファイルを生成すると解釈できる：環境を与えられると望んだ結果を生成するアクションだ。そのようなアクションを実行するためのヘルパー関数\code{runReader}がある：

\src{snippet07}
環境の異なる値に対して異なる結果を生成するかもしれない。

\code{Reader}を返す関数も\code{Reader}アクション自体も純粋であることに注意してほしい。

\code{Reader}モナドのバインドを実装するには、まず環境\code{e}を受け取り\code{b}を生成する関数を作らなければならないことに気づく：

\src{snippet08}
ラムダの中でアクション\code{ra}を実行して\code{a}を生成できる：

\src{snippet09}
次に\code{a}を継続\code{k}に渡して新しいアクション\code{rb}を得られる：

\src{snippet10}
最後に、アクション\code{rb}を環境\code{e}で実行できる：

\src{snippet11}
\code{return}を実装するには、環境を無視して変更されていない値を返すアクションを作る。

これらをまとめて、いくつかの単純化を施すと、次の定義が得られる：

\src{snippet12}

\subsection{書き込み専用の状態}

これは最初のロギングの例そのものだ。装飾は\code{Writer}関手によって与えられる：

\src{snippet13}
完全を期すために、データコンストラクターを展開するだけのトリビアルなヘルパー\code{runWriter}もある：

\src{snippet14}
以前見たように、\code{Writer}を合成可能にするために、\code{w}はモノイドでなければならない。バインド演算子を使って書かれた\code{Writer}のモナドインスタンスは以下の通りだ：

\src{snippet15}

\subsection{状態}

状態への読み書きアクセスを持つ関数は、\code{Reader}と\code{Writer}の装飾を組み合わせる。それらを、状態を追加引数として取り、値と状態のペアを結果として生成する純粋な関数と考えることができる：\code{(a, s) -> (b, s)}。カリー化すると、\code{State}関手に抽象化された装飾を持つKleisli矢の形式\code{a -> (s -> (b, s))}になる：

\src{snippet16}
ここでもKleisli矢をアクションを返すものとして見ることができ、ヘルパー関数で実行される：

\src{snippet17}
異なる初期状態は異なる結果だけでなく、異なる最終状態も生成するかもしれない。

\code{State}モナドのバインドの実装は\code{Reader}モナドのものと非常に似ているが、各ステップで正しい状態を渡すことに注意が必要だ：

\src{snippet18}
以下が完全なインスタンスだ：

\src{snippet19}
状態を操作するために使える二つのヘルパーKleisli矢もある。一つは検査のために状態を取得する：

\src{snippet20}
もう一つは状態を完全に新しい状態に置き換える：

\src{snippet21}

\subsection{例外}

例外を投げる命令型関数は実際には部分関数だ――引数のある値に対して定義されていない関数だ。純粋な全関数を使って例外を実装する最も単純な方法は、\code{Maybe}関手を使う。部分関数は全関数に拡張され、意味をなす場合は\code{Just a}を、そうでない場合は\code{Nothing}を返す。失敗の原因に関する情報も返したい場合は、代わりに\code{Either}関手を使うことができる（最初の型を\code{String}などに固定して）。

以下は\code{Maybe}の\code{Monad}インスタンスだ：

\src{snippet22}
\code{Maybe}のモナド的合成は、エラーが検出されたときに（継続\code{k}が呼ばれることなく）計算を正しく短絡させることに注意してほしい。それが例外に期待する動作だ。

\subsection{継続}

就職面接の後に経験するかもしれない「こちらから連絡します」という状況だ。直接の答えをもらう代わりに、結果と共に呼び出されるハンドラー、つまり関数を提供することになる。このプログラミングスタイルは、結果が呼び出し時点では不明な場合、たとえば別のスレッドで評価されていたりリモートのウェブサイトから配信されたりする場合に特に有用だ。この場合のKleisli矢は、「計算の残り」を表すハンドラーを受け取る関数を返す：

\src{snippet23}
ハンドラー\code{a -> r}は最終的に呼び出されると型\code{r}の結果を生成し、この結果が最後に返される。継続は結果型でパラメータ化される（実際には、これはしばしばある種のステータスインジケーターだ）。

Kleisli矢が返すアクションを実行するためのヘルパー関数もある。ハンドラーを受け取り、継続に渡す：

\src{snippet24}
継続の合成は非常に難しいことで有名だ。そのためモナド、特に\code{do}記法を通じた処理は極めて有利だ。

バインドの実装を考えてみよう。まず、簡略化したシグネチャを見てみよう：

\src{snippet25}
目標はハンドラー\code{(b -> r)}を受け取り結果\code{r}を生成する関数を作ることだ。それが出発点だ：

\src{snippet26}
ラムダの内部では、計算の残りを表す適切なハンドラーを使って関数\code{ka}を呼び出したい。このハンドラーをラムダとして実装する：

\src{snippet27}
この場合、計算の残りはまず\code{kab}を\code{a}で呼び出し、次に結果のアクション\code{kb}に\code{hb}を渡すことだ：

\src{snippet28}
見てわかるように、継続は内側から外側に向かって合成される。最終的なハンドラー\code{hb}は計算の最も内側の層から呼び出される。完全なインスタンスは以下の通りだ：

\src{snippet29}

\subsection{対話型入力}

これは最も扱いにくい問題であり、多くの混乱の源だ。明らかに、\code{getChar}のような関数がキーボードで入力された文字を返すとしたら、純粋ではありえない。しかし、コンテナの内側に文字を返すとしたらどうだろう？そのコンテナから文字を取り出す方法がない限り、その関数は純粋だと主張できる。\code{getChar}を呼び出すたびに、まったく同じコンテナを返すだろう。概念的には、このコンテナはすべての可能な文字の重ね合わせを含む。

量子力学に親しんでいるなら、この類比を理解するのに問題はないはずだ。箱の中のシュレーディンガーの猫と同じだ――ただし、箱を開けたり中を覗いたりする方法がない。箱は特別な組み込みの\code{IO}関手を使って定義される。この例では、\code{getChar}はKleisli矢として宣言できる：

\src{snippet30}
（実際には、ユニット型からの関数は戻り値型の値を選ぶことと等価なので、\code{getChar}の宣言は\code{getChar :: IO Char}に単純化される。）

関手として、\code{IO}は\code{fmap}を使ってその内容を操作できる。また関手として、文字だけでなく任意の型の内容を保存できる。このアプローチの真の有用性は、HaskellでIOがモナドであることを考慮すると明らかになる。これは、\code{IO}オブジェクトを生成するKleisli矢を合成できることを意味する。

Kleisli合成によって\code{IO}オブジェクトの内容を覗けると思うかもしれない（量子の類比を続けるなら「波動関数を崩壊させる」ことになる）。実際、\code{getChar}を、文字を受け取ってたとえば整数に変換する別のKleisli矢と合成できる。ただし、この二番目のKleisli矢はこの整数を\code{(IO\ Int)}として返すことしかできない。また、すべての可能な整数の重ね合わせが生じる。そして以降同様だ。シュレーディンガーの猫が袋から出ることはない。一度\code{IO}モナドの中に入ったら、抜け出す方法はない。\code{IO}モナドには\code{runState}や\code{runReader}に相当するものはない。\code{runIO}は存在しない！

では、Kleisli矢の結果である\code{IO}オブジェクトを別のKleisli矢と合成すること以外に何ができるのか？それを\code{main}から返せる。Haskellでは、\code{main}は次のシグネチャを持つ：

\src{snippet31}
そしてKleisli矢として考えることもできる：

\src{snippet32}
その観点から、Haskellプログラムは\code{IO}モナドにおける一つの大きなKleisli矢だ。モナド的合成を使って、より小さなKleisli矢から合成できる。結果の\code{IO}オブジェクト（\code{IO}アクションとも呼ばれる）で何かをするのはランタイムシステムに任せられる。

矢自体は純粋な関数であることに注意してほしい――すべてが純粋な関数だ。汚い仕事はシステムに委ねられている。\code{main}から返された\code{IO}アクションを最終的に実行するとき、ランタイムはユーザー入力の読み取り、ファイルの変更、うるさいメッセージの表示、ディスクのフォーマットなどのあらゆる厄介なことを行う。Haskellプログラムは手を汚すことはない（\code{unsafePerformIO}を呼び出す場合を除くが、それは別の話だ）。

もちろん、Haskellは遅延なので、\code{main}はほぼ即座に返り、汚い仕事はすぐに始まる。純粋な計算の結果が要求され、必要に応じて評価されるのは\code{IO}アクションの実行中だ。したがって実際には、プログラムの実行は純粋な（Haskell）コードと汚い（システム）コードが交互に現れる。

\code{IO}モナドには、さらに奇妙だが数学的モデルとして完璧に意味をなす別の解釈がある。それはプログラム内のオブジェクトとして宇宙全体を扱う。概念的に、命令型モデルは宇宙を外部のグローバルオブジェクトとして扱うため、I/Oを実行するプロシージャはそのオブジェクトと相互作用することで副作用を持つことに注意してほしい。それらは宇宙の状態を読み取り、変更することができる。

関数型プログラミングで状態をどう扱うかはすでに知っている――状態モナドを使う。しかし単純な状態とは異なり、宇宙の状態は標準的なデータ構造を使って容易に記述できない。しかし、直接それと相互作用しない限り、そうする必要はない。型\code{RealWorld}が存在し、何らかの宇宙的なエンジニアリングの奇跡によって、ランタイムがこの型のオブジェクトを提供できると仮定するだけで十分だ。\code{IO}アクションはただの関数だ：

\src{snippet33}
または\code{State}モナドの用語では：

\src{snippet34}
ただし、\code{IO}モナドの\code{>=>}と\code{return}は言語に組み込まれていなければならない。

\subsection{対話型出力}

同じ\code{IO}モナドが対話型出力をカプセル化するために使われる。\code{RealWorld}はすべての出力デバイスを含むと想定される。なぜHaskellから出力関数を呼び出して何もしないふりをするだけではいけないのか疑問に思うかもしれない。たとえば、なぜ次のようにするのか：

\src{snippet35}
単純なものではなく：

\src{snippet36}
二つの理由がある：Haskellは遅延なので、出力――ここではユニットオブジェクト――が何にも使われない関数を呼び出すことはない。そして、遅延でなくても、そのような呼び出しの順序を自由に変えて出力を乱してしまう可能性がある。Haskellで二つの関数の順次実行を強制する唯一の方法はデータ依存性を通じてだ。一方の関数の入力は他方の出力に依存しなければならない。\code{IO}アクション間で\code{RealWorld}が渡されることで順序付けが強制される。

概念的には、このプログラムでは：

\src{snippet37}
「World!」を表示するアクションは、「Hello 」がすでに画面に表示されている宇宙を入力として受け取る。そして「Hello World!」が画面に表示された新しい宇宙を出力する。

\section{まとめ}

もちろん、モナド的プログラミングの表面をかじっただけだ。モナドは、命令型プログラミングで副作用で行われることを純粋な関数で達成するだけでなく、高度な制御と型安全性でそれを行う。しかし欠点もある。モナドについての主な不満は、互いに簡単に合成できないことだ。確かに、モナドトランスフォーマーライブラリを使って基本的なモナドのほとんどを組み合わせることができる。状態と例外を組み合わせるモナドスタックを作るのは比較的簡単だが、任意のモナドを積み重ねる公式は存在しない。
